\documentclass[12pt, twoside]{article}
%\documentclass[12pt, twoside]{article}
\usepackage[letterpaper, margin=1in, headsep=0.2in]{geometry}
\setlength{\headheight}{0.6in}
%\usepackage[english]{babel}
\usepackage[utf8]{inputenc}
\usepackage{microtype}
\usepackage{amsmath}
\usepackage{amssymb}
%\usepackage{amsfonts}
\usepackage[nomessages]{fp} %\FPeval{\var-name}{2*sin(pi/6)}
\usepackage{siunitx} %units in math. eg 20\milli\meter
\usepackage{yhmath} % for arcs, overparenth command
\usepackage{tikz} %graphics
\usetikzlibrary{quotes, angles, arrows, arrows.meta}
\usepackage{pgfplots}
\usepgfplotslibrary{fillbetween}
\pgfplotsset{compat=1.16}
\usepackage{graphicx} %consider setting \graphicspath{{images/}}
\usepackage{parskip} %no paragraph indent
\usepackage{enumitem}
\usepackage{multicol}
\usepackage{venndiagram}

\usepackage{fancyhdr}
\pagestyle{fancy}
\fancyhf{}
\renewcommand{\headrulewidth}{0pt} % disable the underline of the header
\raggedbottom
\hfuzz=2mm %suppresses overfull box warnings

\usepackage{hyperref}

\fancyhead[LO]{La Scuola d'Italia / Dr.\ Huson / IB Math 12 \\* 19 October 2026}
\fancyhead[RO]{\fbox{\begin{minipage}{6cm}\footnotesize Nickname:\\[0.5cm]\end{minipage}}}
\fancyfoot[C]{\thepage}

\begin{document}
\subsubsection*{2.1 Classwork: Right-Angled Triangle Trigonometry}

Set your calculator to degree mode. Give answers correct to three
significant figures.

\begin{enumerate}[itemsep=0.15cm]

%--- Problem 1: Name the sides; sin, cos, tan as fractions; angle from a ratio; side from sin; hypotenuse from cos [11 marks] ---
%--- Source: Original (freshly authored, AA SL 3.2); Haese AA SL Ch. 8 ---
\item[1.]
For an acute angle $\theta$ in a right-angled triangle,
$$\sin\theta = \frac{\text{opposite}}{\text{hypotenuse}}, \qquad
\cos\theta = \frac{\text{adjacent}}{\text{hypotenuse}}, \qquad
\tan\theta = \frac{\text{opposite}}{\text{adjacent}}.$$

\begin{center}
\begin{tikzpicture}[scale=0.2]
  \coordinate (A) at (0,0);
  \coordinate (C) at (8,0);
  \coordinate (B) at (8,15);
  \draw[thick] (A) -- (C) -- (B) -- cycle;
  \draw pic[draw, angle radius=0.25cm] {right angle = A--C--B};
  \draw pic["$\theta$", draw, angle radius=0.6cm, angle eccentricity=1.6] {angle = C--A--B};
  \node[below left] at (A) {$A$};
  \node[below right] at (C) {$C$};
  \node[above right] at (B) {$B$};
  \node[below] at (4,0) {$8$};
  \node[right] at (8,7.5) {$15$};
  \node[above left] at (4,7.5) {$17$};
\end{tikzpicture}
\hfill
\begin{tikzpicture}[scale=0.3]
  \coordinate (P) at (0,0);
  \coordinate (Q) at (9.83,0);
  \coordinate (R) at (9.83,6.88);
  \draw[thick] (P) -- (Q) -- (R) -- cycle;
  \draw pic[draw, angle radius=0.25cm] {right angle = P--Q--R};
  \draw pic["$35^\circ$", draw, angle radius=0.7cm, angle eccentricity=1.7] {angle = Q--P--R};
  \node[below left] at (P) {$P$};
  \node[below right] at (Q) {$Q$};
  \node[above right] at (R) {$R$};
  \node[above left] at (4.9,3.44) {$12\,\text{cm}$};
\end{tikzpicture}
\hfill
\begin{tikzpicture}[scale=0.3]
  \coordinate (D) at (0,0);
  \coordinate (E) at (7,0);
  \coordinate (F) at (7,8.34);
  \draw[thick] (D) -- (E) -- (F) -- cycle;
  \draw pic[draw, angle radius=0.25cm] {right angle = D--E--F};
  \draw pic["$50^\circ$", draw, angle radius=0.6cm, angle eccentricity=1.8] {angle = E--D--F};
  \node[below left] at (D) {$D$};
  \node[below right] at (E) {$E$};
  \node[above right] at (F) {$F$};
  \node[below] at (3.5,0) {$7\,\text{cm}$};
  \node[above left] at (3.5,4.17) {$x\,\text{cm}$};
\end{tikzpicture}
\end{center}

\begin{enumerate}[itemsep=0.1cm]
  \item In triangle $ABC$, write down the length of the hypotenuse, the
  side opposite $\theta$, and the side adjacent to $\theta$.
  \item Write down $\sin\theta$, $\cos\theta$ and $\tan\theta$ as fractions.
  \item Find the size of angle $\theta$.
  \item In triangle $PQR$, find the length $QR$.
  \item In triangle $DEF$, find the value of $x$.
\end{enumerate}

\begin{tikzpicture}
  \draw (-0.6,0) rectangle (14.6,9.4);
  \foreach \y in {1,2,...,9} \draw[dotted] (0,\y)--(14.1,\y);
\end{tikzpicture}

\newpage

%--- Problem 2: Ladder (angle from cos, height); tower (elevation, height, hypotenuse, angle of depression) [10 marks] ---
%--- Source: Original (freshly authored, AA SL 3.2, 3.3); Haese AA SL Ch. 8 ---
\item[2.]
A ladder $6\,\text{m}$ long leans against a vertical wall. The foot of the
ladder is on horizontal ground, $1.8\,\text{m}$ from the base of the wall.
It may help to draw a diagram.

\begin{enumerate}[itemsep=0.1cm]
  \item Find the angle between the ladder and the ground.
  \item Find how far up the wall the ladder reaches.
\end{enumerate}

A vertical tower $BT$ stands on horizontal ground. From a point $P$ on the
ground, $40\,\text{m}$ from $B$, the angle of elevation of the top $T$ is
$28^\circ$. A car is parked at point $C$, on the other side of the tower,
$60\,\text{m}$ from $B$.

\begin{center}
\begin{tikzpicture}
  \coordinate (P) at (-3,0);
  \coordinate (B) at (0,0);
  \coordinate (C) at (4.5,0);
  \coordinate (T) at (0,2.4);
  \draw[thick] (P) -- (C);
  \draw[very thick] (B) -- (T);
  \draw (P) -- (T) -- (C);
  \draw[dotted] (T) -- (4.5,2.4);
  \draw pic[draw, angle radius=0.25cm] {right angle = P--B--T};
  \draw pic["$28^\circ$", draw, angle radius=0.9cm, angle eccentricity=1.45] {angle = B--P--T};
  \node[below] at (P) {$P$};
  \node[below] at (B) {$B$};
  \node[below] at (C) {$C$};
  \node[above] at (T) {$T$};
  \node[below] at (-1.5,0) {$40\,\text{m}$};
  \node[below] at (2.25,0) {$60\,\text{m}$};
  \node[right] at (4.6,1.2) {\emph{diagram not to scale}};
\end{tikzpicture}
\end{center}

\begin{enumerate}[itemsep=0.1cm]
  \item[(c)] Find the height of the tower.
  \item[(d)] Find the distance $PT$.
  \item[(e)] Find the angle of depression of the car $C$ from the top of
  the tower, $T$.
\end{enumerate}

\begin{tikzpicture}
  \draw (-0.6,0) rectangle (14.6,10.4);
  \foreach \y in {1,2,...,10} \draw[dotted] (0,\y)--(14.1,\y);
\end{tikzpicture}

\end{enumerate}
\end{document}
